% 12pt and authblk mirror manuscript_annotated.tex, so the title block below
% typesets at exactly the same sizes as the manuscript's.
\documentclass[12pt,a4paper]{article}
\usepackage[margin=2.5cm]{geometry}
\usepackage{authblk}
\usepackage[T1]{fontenc}
\usepackage[utf8]{inputenc}
\usepackage{xcolor}
\usepackage{amsmath,amssymb}
\usepackage{booktabs}
\usepackage{enumitem}
\usepackage{graphicx}
\usepackage{xr}
\externaldocument{manuscript_annotated}
\usepackage[colorlinks=true,linkcolor=black,urlcolor=blue]{hyperref}

% Line references into the annotated manuscript.
% Mechanism: \linelabel{ln:xyz} in manuscript_annotated.tex (lineno package) -> \lnp{ln:xyz} here.
% Requires manuscript_annotated.aux, pulled in via xr/\externaldocument above.
% If manuscript_annotated.aux is missing or stale, \ref silently degrades to "??".
% Flag that loudly instead, so a broken reference cannot slip into a version
% that goes out to the reviewers.
\makeatletter
\newcommand{\lnum}[1]{%
  \@ifundefined{r@#1}%
    {\textcolor{red}{\textbf{[line ??\ -- rebuild manuscript_annotated.tex]}}}%
    {line~\ref{#1}}}
\makeatother
\newcommand{\lnp}[1]{(\lnum{#1})}

\definecolor{revised}{rgb}{0,0,0.75}
\definecolor{todocol}{rgb}{0.75,0.25,0}

% Math symbol macros, mirroring manuscript_annotated.tex
\newcommand{\vectorsym}[1]{\boldsymbol{#1}}
\newcommand{\matrixsym}[1]{\boldsymbol{#1}}

% Marker for comments that are not yet addressed in the manuscript.
\newcommand{\todoitem}[1]{\par\vspace{0.3em}\noindent\textcolor{todocol}{\textbf{[TO BE DONE]} #1}\par}

\setlength{\parindent}{0pt}
\setlength{\parskip}{0.6em}

\newcommand{\comm}[2]{%
  \vspace{1.2em}%
  \noindent\textbf{Comment #1}\par\vspace{0.2em}%
  \textit{#2}\par}

\newenvironment{response}%
  {\begin{list}{}{\setlength{\leftmargin}{1.5em}\setlength{\rightmargin}{0em}}\item[]}%
  {\end{list}}

\newenvironment{revisedtext}%
  {\begin{list}{}{\setlength{\leftmargin}{1.5em}\setlength{\rightmargin}{0em}}\item[]\color{revised}}%
  {\end{list}}

% Title block styled as in manuscript_annotated.tex: \bfseries title, authblk
% authors with superscript affiliation markers, affiliations in small italics.
\title{ \bfseries Response to Reviewers \#1\\[0.4em]
\large Manuscript: ``Working Title of the Paper''}
\author[a]{First Author}
\author[b]{Second Author}
\author[a]{Third Author}

\renewcommand\Affilfont{\small\itshape}
\affil[a]{First Institute, First University, Country}
\affil[b]{Second Institute, Second University, Country}

\date{}

\begin{document}

% The VS Code recipe runs bibtex unconditionally. This document has no
% bibliography, so bibtex would abort on a missing \bibdata/\bibstyle/\citation
% and take the whole build down with it. Writing the three commands into the
% .aux directly keeps bibtex happy. The resulting .bbl is never read, so no
% reference list is typeset.
\makeatletter
\immediate\write\@auxout{\string\citation{firstReference2024}}
\immediate\write\@auxout{\string\bibstyle{plain}}
\immediate\write\@auxout{\string\bibdata{../references}}
\makeatother

\maketitle

% Conventions used in this document:
%   \comm{n}{quoted reviewer comment}         -> the comment, in italics
%   \begin{response} ... \end{response}       -> our answer
%   \begin{revisedtext} ... \end{revisedtext} -> the revised manuscript text, in blue
%   \lnp{ln:label}                            -> (line 123) into the annotated manuscript
%   \todoitem{...}                            -> INTERNAL marker, remove before submission

% Each entry is set in the style it describes. Kept as a plain list rather than
% the \comm / response / revisedtext macros, which would drag their own headers
% and vertical spacing into the key.
This response to the review comments adheres to the following style code:

\begin{itemize}[leftmargin=1.5em, topsep=0.4em, itemsep=0.2em, parsep=0pt]
  \item \textit{Reviewer comments are quoted verbatim, in italics.}
  \item Our answers follow in upright text.
  \item \textcolor{revised}{Revised manuscript text is quoted in blue.}
\end{itemize}

Line numbers refer to the annotated version of the revised manuscript.
\pagebreak

%%%%%%%%%%%%%%%%%%%%%%%%%%%%%%%%%%%%%%%%%%%%%%%%%%%%%%%%%%%%%%%%%%%%%%%%%%%%%%
\section*{Response to Editor}
%%%%%%%%%%%%%%%%%%%%%%%%%%%%%%%%%%%%%%%%%%%%%%%%%%%%%%%%%%%%%%%%%%%%%%%%%%%%%%

\comm{E.1}{Quote the editor's request here.}

\begin{response}
Answer it, and say where in the manuscript the change landed \lnp{ln:abstract}.
\end{response}

%%%%%%%%%%%%%%%%%%%%%%%%%%%%%%%%%%%%%%%%%%%%%%%%%%%%%%%%%%%%%%%%%%%%%%%%%%%%%%
\section*{Reviewer 1}
%%%%%%%%%%%%%%%%%%%%%%%%%%%%%%%%%%%%%%%%%%%%%%%%%%%%%%%%%%%%%%%%%%%%%%%%%%%%%%

\comm{1.1}{Quote the first comment verbatim, so the reviewer can recognise it.}

\begin{response}
Answer the comment directly. If the manuscript changed, point at the line and
quote the new text.
\end{response}

\begin{revisedtext}
The revised sentence, copied from the manuscript \lnp{ln:gap}.
\end{revisedtext}

\comm{1.2}{Quote the second comment.}

\begin{response}
Where a comment is declined, say so plainly and give the reason rather than
deferring it.
\end{response}

%%%%%%%%%%%%%%%%%%%%%%%%%%%%%%%%%%%%%%%%%%%%%%%%%%%%%%%%%%%%%%%%%%%%%%%%%%%%%%
\section*{Reviewer 2}
%%%%%%%%%%%%%%%%%%%%%%%%%%%%%%%%%%%%%%%%%%%%%%%%%%%%%%%%%%%%%%%%%%%%%%%%%%%%%%

\comm{2.1}{Quote the comment.}

\begin{response}
Answer it \lnp{ln:methods}.
\end{response}

\comm{2.2}{Quote the comment.}

\begin{response}
Answer it \lnp{ln:results}.
\end{response}

\end{document}
